\documentclass[letterpaper]{article}

\usepackage[letterpaper,margin=1in]{geometry}

% Compile from the repository root so this path resolves to the pinned WSP
% submodule.
\input{wsp/testing/test-case-library.tex}

\begin{document}
\def\TCID{TC-0004}
\def\TCTitle{Library DLL, version resource and registration}
\def\TCPurpose{Verify that \texttt{mswinsck.ocx} exports the COM entry
points, carries the version and target it was built from, makes a control
through its class factory, and registers and unregisters itself under WFC's
own and the Microsoft identifiers.}
\def\TCPriority{\TCPriorityHigh}

\def\TCRequirementRef{REQ-0002}
\def\TCDesignRef{Not applicable.}
\def\TCFeatureRef{\texttt{src/registration.cpp};
\texttt{src/wfcmswinsck.def}; \texttt{version.rc.in};
\texttt{tests/mswinsck\_dll\_test.c}; \texttt{tests/dll\_check.h}.}
\def\TCTestTechnique{Black-box testing of the DLL from a C program that uses
it as any other program would; round-trip testing of the registration.}

\def\TCPreconditions{Windows; the library is configured with
\texttt{WFC\_MSWINSCK\_SHARED=ON} and the DLL is built; the registry of the
current user does not already hold \texttt{Software\textbackslash WFC%
\textbackslash Libraries\textbackslash mswinsck}.}
\def\TCEnvironment{Windows; registered only for the shared build. The test
holds the CTest resource lock \texttt{registry} so that two registration
tests do not run together. The program writes only under the current user
(\texttt{WFC\_REGISTER\_SCOPE=user}).}
\def\TCAssumptions{The registry of the current user is writable.}

\def\TCInputData{The path of the DLL, the contents of \texttt{VERSION} and of
\texttt{TARGET}; the class \texttt{WFC.MSWinsockLib.Winsock} with its class
id, and the Microsoft ProgID \texttt{MSWinsockLib.Winsock.1} with its class
id.}
\def\TCInitialState{The library is not registered for the current user.}

\def\TCProcedure{\begin{enumerate}
  \item Run \texttt{mswinsck\_dll\_test} (CTest \texttt{TC-0004-dll}, in WFC \texttt{wfc-mswinsck-TC-0004-dll}) with
    the DLL, the version and the target as arguments.
  \item Load the DLL; check that \texttt{DllRegisterServer},
    \texttt{DllUnregisterServer}, \texttt{DllGetClassObject} and
    \texttt{DllCanUnloadNow} exist, that \texttt{DllCanUnloadNow} answers
    \texttt{S\_OK} and that \texttt{DllGetClassObject} answers
    \texttt{CLASS\_E\_CLASSNOTAVAILABLE} for an unknown class id.
  \item Read the version resource: the file version equals the version given
    and the \texttt{Comments} string names the target.
  \item Register with \texttt{WFC\_REGISTER\_MICROSOFT\_IDS=0}; check the
    library key (version, \texttt{Target}, \texttt{Server}) and the ProgID
    and class id keys; unregister and check that nothing is left.
  \item If the Microsoft class id is not in the registry of the current
    user, repeat step 4 with \texttt{WFC\_REGISTER\_MICROSOFT\_IDS=1}; the
    Microsoft ProgID and class id are served by the DLL and removed again.
  \item Create the control through the class factory: it keeps a Long written
    to \texttt{LocalPort}, accepts the method \texttt{Close}, does not know an
    unknown name, refuses aggregation, and \texttt{DllCanUnloadNow} answers
    \texttt{S\_OK} after the release.
\end{enumerate}}
\def\TCExpectedResult{The program exits with status 0 and prints
\texttt{mswinsck: exports, version \ldots{} and registration are right} and
the line for the created control. Exit status 77 means the registry already
held the library, and CTest then reports the test as skipped.}
\def\TCPostConditions{The registry of the current user holds no key of the
library; the DLL stays in the build folder.}

\TCTemplate

\end{document}